\section{Secondary theoretical statements}
\label{app:secondary-theory-statements}

This section preserves the auxiliary static, recovery, probe, and optimization
statements referenced by the main text.  They are separated from the central
optimizer-response theorem because the static change of basis is SSMF prior
art and the positive results use standard additional assumptions.  Complete
proofs and failure cases follow in Appendix~\ref{app:theory}.

\subsection{Static router-coordinate consequences}

The rows $b_1,\ldots,b_K$ are affinely independent when
$[B,\ones]$ has row rank $K$.

\begin{proposition}[Local form of the SSMF gauge]
\label{thm:continuous_gauge}
Assume $K\ge2$, $A$ is strictly positive and has column rank $K$, and the
rows of $B$ are affinely independent.  For every nonzero
$H\in\RR^{K\times K}$ satisfying $H\ones=0$, and every sufficiently small
$t\ne0$,
\begin{equation}
 M_t=I+tH,\qquad A_t=AM_t,\qquad B_t=M_t^{-1}B
 \label{eq:gauge_path}
\end{equation}
is a valid factorization with $A_t\in\simplex_K^L$, affinely independent
bases, and $A_tB_t=AB$.  It is not permutation equivalent to $(A,B)$.  The
local equivalence class therefore has at least $K(K-1)$ continuous
directions.
\end{proposition}

The condition $H\ones=0$ preserves row sums; positivity and invertibility
survive a small perturbation; and
$[B_t,\ones]=M_t^{-1}[B,\ones]$ preserves affine independence.  Full column
rank of $A$ rules out a hidden permutation.

For $0<s<1$, the gauge
\begin{equation}
 M_s=(1-s)I+s\ones\ones^\top/K
 \label{eq:uniformizing_gauge}
\end{equation}
moves every router row toward uniform and strictly increases its entropy
unless it was already uniform, while $M_s^{-1}B$ keeps $\Theta$ fixed.
This uniformizing path preserves argmax assignments. They are not
invariant under other feasible gauges, however. For example,
\begin{equation}
 A=\begin{bmatrix}.1&.9\\.3&.7\\.4&.6\\.9&.1\end{bmatrix},\qquad
 M=\begin{bmatrix}.9&.1\\.4&.6\end{bmatrix}
 \label{eq:argmax_example}
\end{equation}
gives argmax labels $(2,2,2,1)$ before and $(2,1,1,1)$ after multiplication
by $M$.  Yet $(AM)(M^{-1}B)=AB$ for every $B$.

This ambiguity does not erase every partition statistic.  Because $M$ is
invertible, $a_i=a_j$ iff $a_iM=a_jM$; affine independence of the rows of
$B$ already makes this equivalent to $\theta_i=\theta_j$.  Indeed,
$qB=0$ and $q\ones=0$ imply $q=0$ when $[B,\ones]$ has row rank $K$.

\begin{proposition}[End-to-end impossibility]
\label{thm:end_to_end_impossibility}
Suppose an architecture and its observation law depend on $(A,B)$ only through
$\Theta=AB$.  If two gauge-equivalent factorizations have different router
summaries, no estimator based on end-to-end input--output observations can
recover the correct summary in both worlds.  This remains true with oracle
access to the predictor on every input.
\end{proposition}

This result concerns router coordinates; it neither equates effective-weight
equality with functional equivalence nor asserts that either equals a planted
generative graph.  The latter is independently impossible under composition:
a two-layer linear teacher with $W_1=W_2=I$ (tied) and one with
$W_1=2I,\ W_2=I/2$ (untied) both implement $x\mapsto x$.

\subsection{Positive recovery under additional information}

Call a factorization \emph{separable} if every basis has an anchor layer
$i_j$ with $a_{i_j}=e_j$.

\begin{proposition}[Standard separable recovery]
\label{thm:anchor_recovery}
If two $K$-separable factorizations of the same $\Theta$ have affinely
independent bases, they are permutation equivalent.  Consequently, if routing
is one-hot with discrete label $c_i\in\{1,\ldots,K\}$, every basis is used,
and the bases are affinely independent, then
\begin{equation}
 c_i=c_{i'}\quad\Longleftrightarrow\quad\theta_i=\theta_{i'},
 \label{eq:onehot_partition}
\end{equation}
so the sharing partition is recoverable up to basis relabeling.
\end{proposition}

This is standard separable-SSMF geometry specialized to layer rows, not a new
simplex-identification result~\citep{abdolali2021simplex}.  Separability makes
$\operatorname{conv}\{\theta_i\}$ equal the simplex whose vertices are the
bases; vertices are unique, and affine independence gives unique barycentric
coordinates.  The restriction to the separable model class is essential: an
unrestricted soft factorization may place a larger simplex outside the
observed convex hull.

\begin{theorem}[Stable noisy one-hot recovery]
\label{thm:noisy_recovery}
Assume $K\ge2$.  Suppose $c_i\in\{1,\ldots,K\}$ and
$\widehat\theta_i=b_{c_i}+e_i$, $\|e_i\|_2\le\varepsilon$, and
$\gamma=\min_{j\ne k}\|b_j-b_k\|_2>4\varepsilon$.  Connecting two layers
when their observed parameters are at distance at most $2\varepsilon$
recovers the exact one-hot partition.  Each cluster mean estimates its basis
to error at most $\varepsilon$.
\end{theorem}

Parameter access is not always available.  Shared-input probes instead target
the clean functional partition defined in Section~\ref{sec:problem}.  Let
$Y_i(H)$ be the possibly noisy observation of layer $i$'s residual update on
a common probe.  The result below applies only when the distances induced by
that observation law retain a positive gap for the clean target partition;
layer-dependent noise can destroy this condition.  Define
\begin{equation}
 d_{ij}^2=\mathbb E\|Y_i(H)-Y_j(H)\|_2^2,\qquad
 \widehat d_{ij}^2=\frac1n\sum_{s=1}^n
 \|Y_i(H_s)-Y_j(H_s)\|_2^2.
 \label{eq:probe_distances}
\end{equation}

\begin{theorem}[Finite-probe partition recovery]
\label{thm:probe_recovery}
Let $L\ge2$, let $0<\eta<1$, and suppose the target partition has at least
two classes. Let it have population within-class radius
$w=\max_{i\sim_f j}d_{ij}^2$ and between-class separation
$b=\min_{i\not\sim_f j}d_{ij}^2$, with gap $\delta=b-w>0$.  Suppose the
centered squared-distance observations satisfy, uniformly over layer pairs,
the Bernstein tail
\begin{equation}
 \Pr\!\left(\left|\widehat d_{ij}^2-d_{ij}^2\right|>u\right)
 \le2\exp\!\left[-n\min\!\left\{\frac{u^2}{2\nu^2},
                                      \frac{u}{2c}\right\}\right].
 \label{eq:probe_bernstein}
\end{equation}
If, for failure probability $\eta$,
\begin{equation}
 n\ge
 \max\!\left\{\frac{32\nu^2}{\delta^2},\frac{8c}{\delta}\right\}
 \log\frac{2\binom{L}{2}}{\eta},
 \label{eq:probe_sample_complexity}
\end{equation}
then thresholding $\widehat d_{ij}^2$ at $(w+b)/2$ recovers the functional
partition with probability at least $1-\eta$.  The same strict separation
guarantees recovery by single, complete, or average agglomerative linkage when
the number $K_f$ of functional equivalence classes is known.  In general
$K_f$ need not equal the number $K$ of router bases.
\end{theorem}

The result follows by a union bound and certifies equivalence only under the
specified probe distribution; behavior off its support remains unidentified.

\subsection{Optimization and initialization effects}

Let $G=\nabla_\Theta\mathcal L$.

\begin{proposition}[Factorization-induced dynamics]
\label{prop:factor_dynamics}
For Euclidean gradient flow on unconstrained factors, without factor-specific
regularizers,
\begin{equation}
 \dot\Theta=-G(B^\top B)-(AA^\top)G.
 \label{eq:effective_gradient_flow}
\end{equation}
Gauge-equivalent initializations therefore generally have different effective
parameter trajectories despite identical initial $\Theta$, predictions, and
loss.
\end{proposition}

\begin{proposition}[Softmax saturation]
\label{prop:softmax_saturation}
Let $a_i=\operatorname{softmax}(z_i/\tau)$ and suppose
$\max_j a_{ij}=1-\varepsilon$.  With
$g_i=\nabla_{\theta_i}\mathcal L$,
\begin{equation}
 \|\nabla_{z_i}\mathcal L\|_2
 \le \frac{2\varepsilon}{\tau}\|Bg_i\|_2.
 \label{eq:saturation_bound}
\end{equation}
Thus the instantaneous logit gradient is attenuated as the current router row
approaches the simplex boundary.  This pointwise bound alone does not imply
that an initial argmax graph remains fixed over a finite training trajectory;
the loss gradients, basis norm, and probability margin can all change.
\end{proposition}

These results justify testing gauge-equivalent initializations, final--initial
router movement, and recovery under deliberately wrong initial graphs.  They
do not prove that gradient descent prefers a universal smooth, banded, or
low-entropy structure.  Nor do reconstruction loss or end-to-end accuracy by
themselves identify router coordinates, the effective-weight equality
partition, functional equivalence, or a latent generative graph.
